\section{USB 移动 HDD 异盘 repo（G:）十重复现}\label{sec:usb-g}

\begin{table}[htbp]
\centering
\caption{G: USB HDD repo 十次墙钟（Disk 1，与 NVMe 物理分离）}
\label{tab:g-raw}
\small
\begin{tabular}{@{}r rr@{}}
\toprule
Run & $T_i$ (s) & $v_i$ (MB/s) \\
\midrule
 1 & 138.650 & 36.99 \\
 2 & 141.969 & 36.12 \\
 3 & 148.785 & 34.47 \\
 4 & 149.008 & 34.42 \\
 5 & 160.359 & 31.98 \\
 6 & 174.982 & 29.31 \\
 7 & 182.882 & 28.04 \\
 8 & 177.572 & 28.88 \\
 9 & 151.877 & 33.77 \\
10 & 155.055 & 33.08 \\
\midrule
均值 $\bar{T}$ & 158.114 & 32.44 \\
$s$ & 15.418 & — \\
$\mathrm{CV}$ & \multicolumn{2}{c}{9.75\%} \\
极差 & 138.650 -- 182.882 & 28.04 -- 36.99 \\
\bottomrule
\end{tabular}
\end{table}

\begin{figure}[htbp]
\centering
\begin{tikzpicture}
\begin{axis}[
  width=0.92\linewidth,
  height=6.5cm,
  xlabel={Run},
  ylabel={墙钟 $T_i$ (s)},
  xmin=0.5, xmax=10.5,
  ymin=130, ymax=190,
  xtick={1,...,10},
  grid=major,
  grid style={dashed,gray!30},
  legend pos=north west,
  legend style={font=\footnotesize}
]
\addplot[USBColor, thick, mark=*, mark size=2pt]
  coordinates {(1,138.650) (2,141.969) (3,148.785) (4,149.008) (5,160.359)
               (6,174.982) (7,182.882) (8,177.572) (9,151.877) (10,155.055)};
\addlegendentry{G: USB HDD}
\addplot[NVMeColor, dashed, thick] coordinates {(1,19.025) (10,19.025)};
\addlegendentry{E: 均值 19.0\,s}
\end{axis}
\end{tikzpicture}
\caption{G: 十次墙钟序列；虚线为 NVMe E: 对照均值。}
\label{fig:g-timeseries}
\end{figure}

\begin{figure}[htbp]
\centering
\begin{tikzpicture}
\begin{axis}[
  width=0.75\linewidth,
  height=7cm,
  ybar,
  bar width=28pt,
  ymin=0, ymax=300,
  ylabel={$\bar{v}$ (MB/s)},
  symbolic x coords={NVMe E:,USB G:},
  xtick=data,
  nodes near coords,
  grid=major,
  grid style={dashed,gray!30},
  legend style={at={(0.5,-0.18)},anchor=north,legend columns=1,font=\footnotesize},
  enlarge x limits=0.35
]
\addplot[fill=NVMeColor!45,draw=NVMeColor] coordinates {(NVMe E:,269.56)};
\addplot[fill=USBColor!45,draw=USBColor] coordinates {(USB G:,32.44)};
\legend{NVMe 同盘, USB 异盘}
\end{axis}
\end{tikzpicture}
\caption{异盘对照：$\bar{v}$ 阶跃（相对 E: 下降 88\%）。}
\label{fig:cross-disk}
\end{figure}

\begin{theorem}[异盘瓶颈阶跃]
\label{thm:cross-disk}
源在 NVMe 页缓存、repo 迁至 USB HDD（G:）时，
\[
  \vsat^{\mathrm{G}} \approx 32.44\,\text{MB/s}
  \;\ll\;
  \vsat^{\mathrm{E}} \approx 269.56\,\text{MB/s},
\]
且 $\bar{T}^{\mathrm{G}}/\bar{T}^{\mathrm{E}} \approx 158.114/19.025 \approx 8.31$。
\end{theorem}

\begin{proof}
\cref{tab:e-raw,tab:g-raw} 与 \cref{eq:throughput}。
相对吞吐下降 $(1-32.44/269.56)\times100\%\approx 88\%$。
由 \cref{prop:bottleneck}，热缓存下读非瓶颈，$v_{\mathrm{write}}^{\mathrm{USB}}\approx 32$\,MB/s
支配 $\vsat$。
\end{proof}

\begin{remark}[USB 更高 CV]
G: 组 $\mathrm{CV}=9.75\%$，高于 NVMe E: 的 1.24\%，与 USB 总线调度、HDD 寻道及 Run6--7
（174--183\,s）慢点一致，仍满足「饱和平台 + 硬件抖动」，而非算法随机性。
\end{remark}
